% Spanish for Beginners — Week 2 Study Guide (Unit 1 continued)
% Compile: pdflatex unit-1-numeros-nacionalidades.tex

\documentclass[10pt,a4paper]{article}
\usepackage[utf8]{inputenc}
\usepackage[T1]{fontenc}
\usepackage[margin=1.2cm]{geometry}
\usepackage{enumitem}
\setlist{nosep,leftmargin=*}
\usepackage{booktabs}
\usepackage{xcolor}
\usepackage{fancyhdr}
\usepackage{microtype}
\definecolor{w2teal}{RGB}{0,100,100}
\pagestyle{fancy}
\fancyhf{}
\fancyfoot[L]{\small Spanish for Beginners | Week 2}
\fancyfoot[R]{\small\thepage}
\raggedright

\begin{document}
\begin{center}
{\large\bfseries\color{w2teal} Week 2 -- Unit 1 (n\'umeros, nacionalidades, lenguas)}\\[0.2em]
\footnotesize N\'umeros 0--50, nacionalidades, profesiones, 5 lenguas de Espa\~na.
\end{center}\vspace{0.5em}

\section*{\color{w2teal}N\'umeros (0--50)}
0--15: cero, uno, dos, tres, cuatro, cinco, \textbf{seis}, siete, ocho, nueve, diez, once, doce, trece, catorce, quince. \\
16--19: diecis\'eis, diecisiete, dieciocho, diecinueve. \\
20, 30, 40, 50: veinte, treinta, cuarenta, cincuenta. \quad \textit{6 = seis.}

\section*{\color{w2teal}Nacionalidades \& profesiones}
\textbf{Terminaciones:} -o (m) / -a (f): argentino/a, portugu\'es/a. Invariable: belga, canadiense, estadounidense. \\
\textbf{Profesiones:} periodista, arquitecto/a, estudiante, profesor/a, secretaria, cient\'ifica, camarero/a, cocinero/a. \\
\textbf{M\'as:} Colombia $\to$ colombiano/a. Italia $\to$ italiano/a. Alemania $\to$ alem\'an/alemana. Francia $\to$ franc\'es/francesa.

\section*{\color{w2teal}5 lenguas oficiales de Espa\~na}
Espa\~nol, euskera, catal\'an, aran\'es, gallego. \\
Rom\'anicas: espa\~nol, catal\'an, gallego, aran\'es. Euskera: origen incierto. \\
Euskera: Pa\'is Vasco, Navarra. Catal\'an: Catalu\~na, Valencia, Baleares. Aran\'es: Val d'Aran.
\end{document}
